% !TeX program = xelatex

% ============================================================ %
% PREAMBLE (antes preamble.tex, fusionado aquí)                %
% ============================================================ %

% ------------------------------------------------------------ %
% TIPO Y PAQUETES                                              %
% ------------------------------------------------------------ %

\documentclass[11pt, a4paper]{report}
\usepackage[spanish, es-tabla, activeacute]{babel}
% \usepackage[utf8]{inputenc}
\usepackage[pagestyles]{titlesec}
\usepackage[hidelinks]{hyperref}
\usepackage{graphicx}
\usepackage{geometry}
\usepackage{indentfirst}
\usepackage{xcolor}
\usepackage{verbatim}
\usepackage{keyval}
\usepackage{etoolbox}
\usepackage{fontspec}
\usepackage{titling}
\usepackage{tikz}
\usepackage{caption}
\usepackage{chngcntr}
\usepackage[bottom]{footmisc}
\usetikzlibrary{positioning}

\definecolor{eus-color}{gray}{0.4}
\definecolor{eng-color}{gray}{0.5}
\definecolor{ehu-blue}{RGB}{27,74,120}

\newcommand*{\noaddvspace}{\renewcommand*{\addvspace}[1]{}}
\addtocontents{lof}{\protect\noaddvspace}
\addtocontents{lot}{\protect\noaddvspace}

% ------------------------------------------------------------ %
% NUMERACIÓN CONTINUA DE FIGURAS Y TABLAS                      %
% ------------------------------------------------------------ %

\counterwithout{figure}{chapter}
\counterwithout{table}{chapter}

% ------------------------------------------------------------ %
% PARTICIÓN DE PALABRAS AL FINAL DE LÍNEA                      %
% ------------------------------------------------------------ %

\pretolerance=2000
\tolerance=3000

% ------------------------------------------------------------ %
% FUENTES DEL DOCUMENTO                                        %
% ------------------------------------------------------------ %

% Tipografía Computer Modern Roman (vía CM-Unicode), la clásica de LaTeX,
% en vez de la sans-serif corporativa EHUSans, para un aspecto más formal
% y en línea con el de otras memorias de TFG (p.ej. la de Oiher).
\setmainfont[
	Path          = C:/Users/peiol/AppData/Local/Programs/MiKTeX/fonts/opentype/public/cm-unicode/,
	Ligatures     = TeX,
	UprightFont   = cmunrm.otf,
	BoldFont      = cmunbx.otf,
	ItalicFont    = cmunti.otf,
	BoldItalicFont= cmunbi.otf
]{CMU Serif}

\newfontfamily{\sansthick}[
 Path          = fonts/,
 BoldFont      = EHUSans-Bold.otf,
 UprightFont   = EHUSans-Regular.otf
 ]{EHUSans}

 \newfontfamily{\sansblack}[
 Path          = fonts/,
 UprightFont   = EHUSans-Black.otf
 ]{EHUSans}

 \newfontfamily{\serif}[
 Path          = fonts/,
 Ligatures     = TeX,
 UprightFont   = EHUSerif-Light.otf,
 BoldItalicFont= EHUSerif-BoldItalic.otf,
 BoldFont      = EHUSerif-Bold.otf,
 ItalicFont    = EHUSerif-LightItalic.otf
 ]{EHUSerif}

 \newfontfamily{\serifthick}[
 Path          = fonts/,
 BoldFont      = EHUSerif-Bold.otf,
 UprightFont   = EHUSerif-Regular.otf
 ]{EHUSerif}

 \newfontfamily{\serifblack}[
 Path          = fonts/,
 UprightFont   = EHUSerif-Black.otf
 ]{EHUSerif}

\newfontfamily{\calibrifont}{Calibri}

% ------------------------------------------------------------ %
% ETIQUETAS EN LAS TABLAS Y FIGURAS							   %
% ------------------------------------------------------------ %

\DeclareCaptionFont{ehu}{\fontsize{9}{10}\mdseries}
\usepackage[font={ehu}, labelfont={bf}]{caption}

% ------------------------------------------------------------ %
% AÑADIR CUARTO NIVEL DE TÍTULO                                %
% ------------------------------------------------------------ %

\setcounter{secnumdepth}{4}
\titleformat{\paragraph}
{\normalfont\normalsize\bfseries}{\theparagraph}{1em}{}
\titlespacing*{\paragraph}
{0pt}{3.25ex plus 1ex minus .2ex}{1.5ex plus .2ex}

% ------------------------------------------------------------ %
% ELIMINAR CABECERA DE CAPÍTULOS                               %
% ------------------------------------------------------------ %

\titleformat{\chapter}%
  {\normalfont\bfseries\Huge}{\thechapter.}{10pt}{}
\newpagestyle{tfg}{
  \setheadrule{0.4pt}
  \sethead{Trabajo Fin de Grado}{}{}
  \setfootrule{0.4pt}
  \setfoot{UPV/EHU}{}{\thepage}
}

% ------------------------------------------------------------ %
% MÁRGENES Y DISTANCIA ENTRE PÁRRAFOS						   %
% ------------------------------------------------------------ %

\geometry{top=2.5cm,left=3cm,right=3cm,bottom=2.5cm}
\setlength{\parskip}{8pt}

% ------------------------------------------------------------ %
% TÍTULO DEL ÍNDICE                                            %
% ------------------------------------------------------------ %

\addto\captionsspanish{
	\renewcommand{\contentsname}{Índice}
	\renewcommand{\listtablename}{Lista de tablas}
	\renewcommand{\listfigurename}{Lista de figuras}
}

% ------------------------------------------------------------ %
% DIRECTORIO DE IMÁGENES                                       %
% ------------------------------------------------------------ %

\graphicspath{{images/}}

% ------------------------------------------------------------ %
% PORTADA                                                      %
% ------------------------------------------------------------ %

\newcommand{\portada}{
	{\raggedleft\includegraphics[width=8cm]{corp/ehu}}
	\\[1.5\baselineskip]
	\noindent\makebox[0pt][l]{\hspace*{-3cm}\colorbox{ehu-blue}{\parbox{\paperwidth}{\centering\color{white}
		\vspace{0.4em}
		\large{\MakeUppercase{Grado en Ingeniería Informática de Gestión y Sistemas de Información}}
		\\[0.6em]
		\LARGE{\textbf{\MakeUppercase{Trabajo de fin de grado}}}
		\vspace{0.4em}
	}}}
	\\[0.5\baselineskip]
	\begin{center}
		\parbox{13cm}{\centering\LARGE{\calibrifont\bfseries\MakeUppercase{Conversión de actas del Ayuntamiento de Bilbao a Knowledge Graphs (RDF, OWL, SPARQL, SHACL) y vectorial mediante LLMs, enlazado con Wikidata y ofreciendo funcionalidad RAG. Interfaz web: visual, consultas, conversacional}}}
	\end{center}
	\renewcommand{\baselinestretch}{1.25}
	\begin{tikzpicture}[remember picture, overlay]
		\node (C) at ([yshift=-6.95cm, xshift=9.7cm]current page.west) [inner sep=10pt, minimum width=13.34cm, minimum height=64pt, align=left]{};
		\node (D) [right=0.2cm of C.west, anchor=west, align=left]
		{\small{\textbf{Alumno: }Lopez, Sanchez, Peio}\\
			\small{\textbf{Director: }Egaña Aranguren, Mikel}
			% Comentar la siguiente línea si sólo hay una dirección del TFG.
			%\\\small{\textbf{Director/Directora (2): }<apellido 1, apellido 2, nombre>}
			};
		\draw[line width=0.5pt] (C.south west) -- (C.south east);
	\end{tikzpicture}
	\renewcommand{\baselinestretch}{1}
	\begin{tikzpicture}[remember picture, overlay]
		\node (E) at ([yshift=-9.4cm, xshift=8.5cm]current page.west) [inner sep=10pt, minimum width=13.34cm, minimum height=27pt, align=left]{};
		\node (F) [right=1.4cm of E.west, anchor=west]
		{\small{\textbf{Curso: }2026-2027}};
		\draw[line width=0.5pt] (E.south west) -- (E.south east);
	\end{tikzpicture}
	\begin{tikzpicture}[remember picture, overlay]
		\node (G) at ([yshift=-11.6cm, xshift=13.4cm]current page.west) [inner sep=10pt, minimum width=7.93cm, minimum height=27pt, align=left]{};
		\node (H) [left=0.2cm of G.east, anchor=east]
		{\small{\textbf{Fecha: }<XXXX, día, mes, año>}};
		\draw[line width=0.5pt] (G.south west) -- (G.south east);
	\end{tikzpicture}
}

% ============================================================ %
% FIN DEL PREAMBLE                                              %
% ============================================================ %

\begin{document}

	\portada
	\thispagestyle{empty}
	\setcounter{page}{0}
	\pagestyle{tfg}

	% ============================================================ %
	% RESUMEN (antes src/abstract.tex)                              %
	% ============================================================ %

	\chapter*{Resumen \normalfont{{\serifthick Laburpena} Abstract}}
	\addcontentsline{toc}{chapter}{Resumen \normalfont{{\serifthick Laburpena} Abstract}}

	\begin{itshape}
	Este trabajo aborda la conversión de las actas del Pleno del Ayuntamiento de
	Bilbao (2007--2026), publicadas en formato PDF, en dos sistemas de recuperación de
	información consultables en lenguaje natural: un RAG vectorial y un sistema GraphRAG.

	Las actas se obtienen mediante la descarga automatizada del portal de actas del
	Ayuntamiento, y de cada una se extraen las proposiciones individuales del Pleno
	junto con sus metadatos: fecha y sesión plenaria, grupo político o particular
	proponente, tema, entidades mencionadas, resultado de la votación (agregado por
	grupo y, cuando el acta lo recoge de forma explícita, nominal por concejal), y
	página y documento de origen. A partir de estos datos se construye un grafo de
	conocimiento RDF/OWL que modela proposiciones, grupos políticos, concejales y
	temas; un modelo de lenguaje clasifica temáticamente cada proposición y extrae
	las entidades que menciona, incorporando dicha información al grafo.

	Sobre este corpus de casi 3.900 proposiciones se implementan dos sistemas de
	recuperación complementarios: un RAG vectorial, con búsqueda híbrida y
	reordenación de resultados, y un sistema GraphRAG, que traduce la pregunta a una
	consulta SPARQL, la ejecuta sobre el grafo mediante razonamiento automático para
	responder preguntas de agregación, y genera la respuesta a partir de los datos
	obtenidos.

	Ambos sistemas citan la página exacta del documento de origen de cada dato
	mostrado, y se exponen mediante una interfaz web conversacional construida con
	Chainlit. El trabajo incluye una evaluación comparativa sistemática de ambas
	arquitecturas sobre un conjunto de preguntas de referencia con respuesta verificada
	a mano contra las actas originales.

	\textbf{Palabras Clave:} GraphRAG, RAG, LLM, grafo de conocimiento, SPARQL.

	%%%%%%%%%%%%%%%%%%%%%%%%%%%%%%%%%%%

	\textcolor{eus-color}{{\serif [PENDIENTE]}}

	\textcolor{eus-color}{{\serif \textbf{Gako-hitzak:}} GraphRAG, RAG, LLM, ezagutza-grafoa, SPARQL.}

	%%%%%%%%%%%%%%%%%%%%%%%%%%%%%%%%%%%

	\textcolor{eng-color}{{\serif [PENDIENTE]}}

	\textcolor{eng-color}{{\serif \textbf{Keywords:}} GraphRAG, RAG, LLM, knowledge graph, SPARQL.}
	\end{itshape}

	\pagebreak

	\setcounter{tocdepth}{1}
	\tableofcontents

	% ÍNDICES:
	\cleardoublepage
	\phantomsection
	\addcontentsline{toc}{chapter}{\listfigurename}
	\listoffigures % Lista de figuras

	\cleardoublepage
	\phantomsection
	\addcontentsline{toc}{chapter}{\listtablename}
	\listoftables % Lista de tablas

	\cleardoublepage

	% ============================================================ %
	% LISTA DE ACRÓNIMOS (antes src/acronimos.tex)                  %
	% ============================================================ %

	\chapter*{Lista de acrónimos}
	\addcontentsline{toc}{chapter}{Lista de acrónimos}

	\begin{description}
		\item [LLM] Large Language Model
		\item [RAG] Retrieval Augmented Generation
		\item [GraphRAG] Graph Retrieval Augmented Generation
		\item [RDF] Resource Description Framework
		\item [OWL] Web Ontology Language
		\item [SPARQL] SPARQL Protocol and RDF Query Language
		\item [SKOS] Simple Knowledge Organization System
		\item [SHACL] Shapes Constraint Language
		\item [LOD] Linked Open Data
		\item [API] Application Programming Interface
		\item [PDF] Portable Document Format
		\item [OCR] Optical Character Recognition
		\item [TFG] Trabajo Fin de Grado
	\end{description}

	% MEMORIA:

	% ============================================================ %
	% INTRODUCCIÓN (antes src/intro.tex)                            %
	% ============================================================ %

	\chapter{Introducción}

	\section{Descripción del problema}

	El portal de actas del Ayuntamiento de Bilbao publica, para cada sesión del Pleno
	desde 2007 hasta la actualidad, varios documentos PDF distintos: el orden del día,
	extractos, el resumen de la sesión y el acta completa. De estos, el acta es el único
	que recoge de forma íntegra las proposiciones, el debate y el resultado de cada
	votación, por lo que es el documento que utiliza este trabajo; cada acta puede
	superar el centenar de páginas, y el corpus completo reúne cerca de veinte años de
	actas y más de 3.900 proposiciones distintas.

	Dentro de cada acta, además, el contenido relevante --las proposiciones y sus
	votaciones-- se mezcla con una cantidad considerable de contenido puramente
	procedimental: listas de asistentes, fórmulas de quórum, aprobación del acta
	anterior, cuestiones de orden o informes que solo se dan por conocidos sin
	votación. Identificar qué parte de cada acta es efectivamente una proposición
	--y separarla del resto-- es en sí mismo parte del problema a resolver.

	A esto se suma un problema práctico adicional: encontrar información concreta
	--qué votó un grupo político sobre un tema, cuántas proposiciones sobre vivienda
	se presentaron en un año determinado, cómo ha evolucionado un asunto a lo largo de
	los mandatos-- exige leer manualmente decenas de documentos extensos, o como mínimo
	saber de antemano en qué fecha y página buscar. No existe ninguna vía para hacer
	preguntas en lenguaje natural sobre este archivo y obtener una respuesta
	verificable, con cifras y citas a la fuente original.

	\subsection{Acceso a la información municipal}

	El objetivo de este TFG es cerrar esa brecha: construir un sistema que permita consultar
	en castellano, mediante preguntas formuladas de forma natural, el contenido íntegro de
	las actas del Pleno, y que responda con datos verificables --enlazados a la página exacta
	del PDF de origen-- en lugar de con texto generado sin respaldo documental.

	\section{Estructura de la memoria}

	Esta memoria se organiza como sigue. El Capítulo~\ref{ch:contexto} introduce los
	conceptos técnicos en los que se apoya el sistema (modelos de lenguaje, RAG,
	grafos de conocimiento y GraphRAG). El Capítulo~\ref{ch:objetivos} concreta los
	objetivos del trabajo y su alcance real. El Capítulo~\ref{ch:beneficios} discute
	los beneficios del sistema y su relación con los Objetivos de Desarrollo Sostenible.
	Los capítulos siguientes describen el estado del arte, el análisis de alternativas
	tecnológicas, los riesgos del proyecto y, con más detalle, la solución desarrollada:
	los dos sistemas de recuperación de información construidos --uno basado en
	búsqueda vectorial y otro en un grafo de conocimiento RDF/OWL consultado por
	SPARQL-- y la interfaz web conversacional que los expone al usuario. La memoria
	cierra con la planificación seguida, el presupuesto, las conclusiones y los anexos.

	% ============================================================ %
	% CONTEXTO (antes src/contexto.tex)                             %
	% ============================================================ %

	\chapter{Contexto}
	\label{ch:contexto}

	A continuación se presentan los conceptos técnicos en los que se apoya este
	trabajo: los modelos de lenguaje y la generación aumentada por recuperación
	(RAG), los grafos de conocimiento, y GraphRAG como la combinación de ambos. Estos
	conceptos son la base sobre la que se construyen los dos sistemas de recuperación
	de información comparados en este TFG.

	\section{Inteligencia Artificial Generativa}

	La Inteligencia Artificial (IA) Generativa agrupa a los modelos capaces de producir
	contenido nuevo --texto, código, imágenes-- a partir de un patrón aprendido de
	grandes cantidades de datos, en lugar de limitarse a clasificar o predecir sobre
	datos ya existentes. Los modelos de lenguaje son el caso particular relevante para
	este trabajo: reciben texto como entrada y generan texto como salida, condicionados
	por una instrucción (\textit{prompt}).

	\section{Modelos de Lenguaje de Gran Tamaño}

	Los Modelos de Lenguaje de Gran Tamaño (LLM, \textit{Large Language Models}) son
	redes neuronales entrenadas sobre corpus masivos de texto que, dada una secuencia de
	entrada, predicen la continuación más probable. Esta capacidad, combinada con un
	ajuste posterior para seguir instrucciones, permite emplearlos para tareas tan
	diversas como resumir, traducir, generar código o --el uso que se les da en este
	TFG-- extraer información estructurada de texto libre y redactar respuestas a
	partir de datos ya recuperados.

	Un LLM no tiene, por sí mismo, acceso a información posterior a su entrenamiento
	ni a documentos privados como las actas de un ayuntamiento concreto. Tampoco
	garantiza que lo que genera sea correcto: puede producir afirmaciones plausibles
	pero falsas, fenómeno conocido como \textit{alucinación}. Ambas limitaciones son
	las que motivan la siguiente técnica.

	\section{Retrieval Augmented Generation (RAG)}

	Retrieval Augmented Generation (RAG) es una arquitectura que combina un sistema de
	recuperación de información con un LLM: antes de generar la respuesta, se buscan
	en una base documental los fragmentos más relevantes para la pregunta y se
	proporcionan al modelo como contexto. El LLM genera entonces la respuesta
	apoyándose únicamente en esos fragmentos, no en su conocimiento interno.

	Esto tiene dos ventajas directas para este proyecto: permite responder sobre un
	corpus privado y actualizable (las actas del Pleno) sin reentrenar ningún modelo,
	y permite citar la fuente exacta de cada afirmación, ya que se conoce qué
	fragmentos --y por tanto qué documento y página-- sustentan la respuesta.

	La recuperación de los fragmentos relevantes suele apoyarse en \textit{embeddings}:
	representaciones vectoriales del texto en las que fragmentos semánticamente
	similares quedan cerca en el espacio vectorial, lo que permite buscar por
	significado y no solo por coincidencia literal de palabras.

	\section{Grafos de Conocimiento}

	Un grafo de conocimiento representa la información como una red de entidades
	(nodos) y relaciones entre ellas (aristas), en lugar de como texto libre. El
	estándar utilizado en este trabajo es RDF (\textit{Resource Description
	Framework}), que modela cada hecho como una tripleta \textit{sujeto-predicado-objeto}
	(por ejemplo, \texttt{proposición~X --presentadaPor--> Grupo~PP}). Sobre RDF se
	puede definir una ontología con OWL (\textit{Web Ontology Language}) --las clases y
	relaciones válidas del dominio-- y consultar el grafo con el lenguaje SPARQL.

	Una ventaja de este modelo frente al texto libre es el razonamiento automático: un
	motor de inferencia (\textit{reasoner}) puede derivar hechos nuevos que no están
	escritos explícitamente pero se siguen de la ontología --por ejemplo, que una
	proposición sobre "desahucios" trata también, por herencia temática, de
	"vivienda"-- sin que haga falta repetir esa relación en cada dato.

	\section{GraphRAG}

	GraphRAG es la variante de RAG en la que la fuente de recuperación no es un
	almacén de fragmentos de texto sino un grafo de conocimiento consultado mediante
	un lenguaje de consulta estructurado (SPARQL en este trabajo). En lugar de
	recuperar pasajes de texto parecidos a la pregunta, un LLM traduce la pregunta en
	lenguaje natural a una consulta formal, el motor de grafos la ejecuta y devuelve
	filas de datos exactos, que un segundo LLM convierte en una respuesta narrada.

	Esto resulta especialmente adecuado para preguntas de agregación y razonamiento
	multi-salto --"¿qué grupo ha presentado más proposiciones sobre vivienda?",
	"¿cuántas se han aprobado por año?"-- que un sistema de recuperación de texto
	libre no puede responder con precisión, porque exigen contar y cruzar datos de
	todo el corpus, no solo encontrar un pasaje parecido a la pregunta. Es, junto con
	el RAG vectorial tradicional, uno de los dos sistemas comparados en este TFG.

	% ============================================================ %
	% OBJETIVOS Y ALCANCE (antes src/objetivos.tex)                 %
	% ============================================================ %

	\chapter{Objetivos y alcance}
	\label{ch:objetivos}

	A continuación se presentan los objetivos definidos para este TFG, así como el
	alcance del sistema finalmente implementado.

	\section{Objetivos del TFG}

	Se presenta el objetivo general del trabajo y los objetivos específicos que se
	derivan de él.

	\section{Alcance del sistema desarrollado}

	Se describe qué incluye y qué queda fuera del sistema final, y por qué -- p.ej.
	el enlazado con Wikidata, que se valoró y se descartó para esta entrega.

	% ============================================================ %
	% BENEFICIOS Y ODS (antes src/beneficios.tex)                   %
	% ============================================================ %

	\chapter{Beneficios que aporta el trabajo y relación con los Objetivos de Desarrollo Sostenible}
	\label{ch:beneficios}

	A continuación se explica a quién beneficia el sistema desarrollado y cómo se
	relaciona con los Objetivos de Desarrollo Sostenible de la Agenda 2030.

	\section{Beneficios del sistema desarrollado}

	Se identifican los distintos perfiles de usuario a los que beneficia el sistema
	y de qué forma.

	\section{Relación con los Objetivos de Desarrollo Sostenible}

	Se relacionan los ODS pertinentes con el trabajo desarrollado.

	% ============================================================ %
	% ESTADO DEL ARTE (nuevo capítulo, esqueleto)                   %
	% ============================================================ %

	\chapter{Descripción de requerimientos y análisis del estado del arte}
	\label{ch:estadodelarte}

	A continuación se describe el corpus documental de partida, los requisitos que
	debía cumplir el sistema y el estado del arte de las tecnologías empleadas.

	\section{Corpus documental}

	Se describen los documentos que publica el portal de actas del Ayuntamiento,
	cuál de ellos usa este trabajo y por qué, y las limitaciones que tiene como
	fuente (PDF escaneado/OCR, falta de estructura, etc.).

	\section{Requisitos del sistema}

	Se enumeran los requisitos funcionales y no funcionales del sistema final.

	\section{Estado del arte de las tecnologías utilizadas}

	Se revisan trabajos y sistemas similares existentes -- otros RAG/GraphRAG sobre
	actas o documentos administrativos -- y se compara con el TFG de Oiher
	Rodríguez Luengo como referencia más cercana.

	% ============================================================ %
	% ANÁLISIS DE ALTERNATIVAS (antes src/alternativas.tex)         %
	% ============================================================ %

	\chapter{Análisis de alternativas}

	A continuación se comparan las principales decisiones tecnológicas evaluadas
	antes de fijar la arquitectura final del sistema.

	\section{Modelos de lenguaje: local (Ollama) frente a API en la nube (Groq / Gemini)}

	\subsection{Puntos positivos}
	\subsection{Puntos negativos}
	\subsection{Conclusión}

	\section{Almacenamiento del grafo: rdflib/OWL-RL frente a una base de grafos dedicada (p.ej. GraphDB, Neo4j)}

	\subsection{Puntos positivos}
	\subsection{Puntos negativos}
	\subsection{Conclusión}

	\section{Almacén vectorial y recuperación: ChromaDB e híbrida (semántica + literal + temática) frente a solo semántica}

	\subsection{Puntos positivos}
	\subsection{Puntos negativos}
	\subsection{Conclusión}

	% ============================================================ %
	% ANÁLISIS DE RIESGOS (nuevo capítulo, esqueleto)                %
	% ============================================================ %

	\chapter{Análisis de riesgos}
	\label{ch:riesgos}

	A continuación se identifican los principales riesgos del proyecto y las
	medidas adoptadas para prevenirlos o mitigarlos.

	\section{Medidas de prevención y contingencia}

	Se detallan los riesgos identificados (p.ej. dependencia de APIs externas,
	calidad del OCR de las actas antiguas, coste o límite de peticiones a los LLM)
	y cómo se mitigó cada uno.

	% ============================================================ %
	% DESCRIPCIÓN DE LA SOLUCIÓN (antes src/solucion.tex)           %
	% ============================================================ %

	\chapter{Descripción de la solución propuesta}
	\label{ch:solucion}

	A continuación se describe en detalle el sistema desarrollado: su arquitectura,
	el pipeline de datos, el grafo de conocimiento, los dos motores de recuperación
	y la interfaz que los expone al usuario.

	\section{Arquitectura general del sistema}

	Se presenta la visión general de los componentes del sistema y cómo se
	conectan entre sí.

	\section{Pipeline de procesamiento de datos}

	Se describe el proceso completo desde el scraping de las actas hasta el
	enriquecimiento de las proposiciones con un LLM mediante un esquema JSON
	cerrado (ver Capítulo~\ref{ch:contexto}).

	\section{Construcción del grafo de conocimiento}

	Se describe la ontología OWL, el paso de JSON a triples RDF, el razonamiento
	automático (\textit{roll-up} temático) y la validación SHACL del grafo.

	\section{Sistema de recuperación de información}

	Se describen los dos motores de recuperación implementados sobre el mismo
	corpus.

	\subsection{RAG vectorial}
	\subsection{GraphRAG}

	\section{Generación de respuestas}

	Se explica cómo se narra la respuesta a partir de los datos recuperados, y
	cómo se citan las fuentes (página y documento exactos) en ambos sistemas.

	\section{Interfaz web del sistema}

	Se describen los componentes de la interfaz Chainlit y el flujo de uso.

	\section{Flujo completo del sistema}

	\section{Implementación técnica}

	% METODOLOGÍA SEGUIDA EN EL DESARROLLO DEL TRABAJO:

	% ============================================================ %
	% PLANIFICACIÓN / GANTT (antes src/gantt.tex)                   %
	% ============================================================ %

	\chapter{Planificación del proyecto}
	\label{ch:planificacion}

	A continuación se describe la planificación temporal seguida durante el
	desarrollo del trabajo.

	\section{Diagrama EDT}

	\section{Descripción de tareas}

	\subsection{Gestión del proyecto}
	\subsection{Investigación y estudio previo}
	\subsection{Desarrollo del sistema}
	\subsection{Evaluación del sistema}
	\subsection{Documentación}

	\section{Diagrama de Gantt}

	\section{Metodología de desarrollo}

	% ============================================================ %
	% ANÁLISIS DE RESULTADOS (nuevo capítulo, esqueleto)             %
	% ============================================================ %

	\chapter{Análisis de resultados}
	\label{ch:resultados}

	A continuación se presenta la evaluación comparativa de los dos sistemas
	desarrollados: RAG vectorial y GraphRAG.

	\section{Objetivo de la evaluación}

	Se explica qué se quiere medir al comparar ambos sistemas.

	\section{Metodología}

	Se describe la batería de preguntas de referencia y su verificación manual
	contra las actas originales en PDF.

	\section{Métricas utilizadas}

	\section{Experimentos realizados}

	\section{Resultados obtenidos}

	\section{Limitaciones}

	% ASPECTOS ECONÓMICOS:

	% ============================================================ %
	% DESCRIPCIÓN DEL PRESUPUESTO (antes src/presupuesto.tex)       %
	% ============================================================ %

	\chapter{Descripción del presupuesto}

	A continuación se detalla el presupuesto estimado del proyecto.

	\section{Costes en software y hardware}

	\section{Costes por mano de obra}

	% CONCLUSIONES:

	% ============================================================ %
	% CONCLUSIONES (antes src/conclusiones.tex)                     %
	% ============================================================ %

	\chapter{Conclusiones}
	\label{ch:conclusiones}

	A continuación se recogen las conclusiones del trabajo, el grado de
	cumplimiento de los objetivos planteados y las mejoras futuras identificadas.

	\section{Conclusiones sobre el sistema}

	\section{Cumplimiento de objetivos}

	\section{Conclusiones sobre las pruebas realizadas}

	\section{Mejoras futuras}

	Se recogen el enlazado con Wikidata sobre entidades no personales, la
	desambiguación por género de concejales homónimos, y otras mejoras apuntadas
	durante el desarrollo.

	% BIBLIOGRAFÍA:

	% ============================================================ %
	% BIBLIOGRAFÍA (antes src/bibliografia.tex)                     %
	% ============================================================ %

	\cleardoublepage
	\phantomsection
	\addcontentsline{toc}{chapter}{Bibliografía}
	\begin{thebibliography}{9}
		\bibitem{ref-1}
		Autores.
		\textit{Título}.
		Editorial, Año.

		\bibitem{ref-2}
		Autores.
		\textit{Título}.
		Editorial, Año.
		\url{https://www.example.com}
	\end{thebibliography}

	% ANEXOS:

	%	\include{src/normativa}
	%	\include{src/planos}
	%	\include{src/manuales}

\end{document}
